\fsection{Descripción inicial}
El cliente es un centro logístico de comercio electrónico en el que los productos llegan a la zona de preparación
de pedidos amontonados y sin orden dentro de contenedores. Hoy en día el picking se hace a mano: un operario coge
cada pieza del contenedor y la deja en la bandeja de salida. Es un trabajo repetitivo y poco ergonómico, que se
convierte en un cuello de botella en los picos de demanda. Los sistemas robotizados convencionales que ha
valorado no le sirven, porque necesitan que las piezas lleguen ordenadas o en una posición conocida.

En la primera visita el cliente nos enseña los contenedores tal como llegan y nos plantea sus necesidades:
\begin{itemize}
	\item Automatizar la recogida de piezas directamente desde el contenedor, sin tener que ordenarlas ni
	      posicionarlas antes, porque los productos cambian con frecuencia y no quiere utillajes específicos
	      para cada uno.
	\item Que la solución pueda convivir con personas en el mismo espacio de trabajo, sin grandes vallados.
\end{itemize}

Para dar respuesta a estas necesidades se propone una celda robotizada de bin picking: una cámara 3D y un modelo
de inteligencia artificial detectan las piezas del contenedor y calculan el mejor punto de agarre, un robot
colaborativo las recoge con una pinza de vacío y las deposita en la bandeja de colocación, mientras un PLC coordina
todo el ciclo. El modelo de IA es el preentrenado por
Siemens, por lo que su entrenamiento queda fuera del alcance del proyecto; el trabajo se centra en la
integración, la programación y la puesta en marcha del conjunto.
